\begin{acknowledgements}


I am deeply grateful to my advisor, Predrag Radivojac, who has been an excellent teacher and mentor.  The course of my life was irrevocably changed when I took his machine learning class and was then asked to work on a project as an undergraduate researcher. This resulted in me subsequently being asked to continue working with him as a masters student, and then as a PhD student.  Over the past 8 years Pedja has been the catalyst behind the growth I've experienced as a scientist and researcher.

In addition to my advisor I have to thank my committee members, Haixu Tang, Matthew Hahn, and Michael Lynch for their contributions.  As well as their input regarding my dissertation, each of these individuals has contributed to my growth through their teaching.  From Haixu Tang I have learned much about algorithms and probabilistic models.  I have tried to absorb as much as I can from Michael Lynch about population genetics and evolution as possible. Matthew Hahn has not only taught excellent courses on evolution, genomics, and population genetics but has also been directly involved in my research on at least one occasion.

I would like to thank Jim Costello, Brian Eades, and John Colbourne who served as mentors for my undergraduate capstone project that focused on using Hidden Markov Models to identify transposable elements.  I'd also like to thank Michael Lynch and his student at that time, Sarah Schaack, for supervising me for one summer in my rather feeble attempt to assist Sarah in quantifying transposable elements in the \emph{Daphnia} genome.  Not only did I learn a lot about conducting bioinformatics research, the time I spent working in the laboratory served as an informal crash course in population genetics and evolution.

Facuty members Mehmet Dalkilic and Sun Kim both contributed to my education through their teaching and involvement in my qualifying exams.  Mehmet was also a research collaborator on several projects; notably the first paper I published as an undergraduate.

I must acknowledge both Iddo Friedberg and Sean Mooney, with whom I have greatly enjoyed collaborating on the CAFA project. I directly benefitted from the hard work that Iddo and Sean, along with my advisor, put into making CAFA successful.


I have to thank my co-author on the ortholog conjecture paper, Nathan Nehrt, who carried out the brunt of the work for that particular publication. 

I am grateful for my friends and classmates Jose Lugo-Martinez, Vikas Rao Pejaver, Heewook Lee, James Denton, Taylor Raborn, Nathan Taylor, Sam Miller, Nathan Husted and Anoop Mayampurath; with whom conversation has always been intellectually stimulating and entertaining. Their camaraderie has made graduate school easier to get through, and they will be missed dearly when I leave Bloomington.  I would also like to recognize both past and present lab members who have not already been mentioned; Amrita Mohan, Biao Li, Yong Li, Fuxiao Xin, Nils Schimmelman, Xin He, Sujun Li, Arunima Ram, Shuyan Li, Kymberleigh Pagel, Chao Ji, and Yuxiang Jiang for their involvement in my research and for making the Radivojac Laboratory an enjoyable and productive environment to work in.

The staff in the School of Informatics and Computing have been very helpful; especially Dave Cooley, Linda Hostetter, Lynne Mikolon, Bob Konicek, Rob Henderson, T.J. Jones and Patty Reyes-Cooksey .  

I am indebted to all teachers at every level; from kindergarden to graduate school. I received a very high-quality public education that served as the foundation for the rest of my career. I especially have to thank Mark Werking, who was an excellent chemistry and physics teacher.  Much of what I learned in the several years I was his student made understanding the abstract concepts and forces that influence protein behavior at the molecular level easy.  I had very good math teachers in Chris Oliger and Tom Hamilton, who helped instill excellent analytical skills.  My French teacher, Wanetta Cheesman, facilitated my growth culturally, especially through organizing a class trip to England France and Switzerland. In the ``shop classes" of Bill Bunger, Don Sturgeon, and Harold Lumpkin I developed pragmatic problem solving skills that I use every day in my research.  I am especially grateful for the personal attention that Harold Lumpikin gave me as a student. In addition to the aforementioned high school teachers, I would like to thank Leah Savion; who in addition to being an excellent teacher, guided me towards informatics as an undergraduate.  Dennis Groth has influenced not only me, but hundreds of students who have come through the School of Informatics and Computing as well.  At the graduate level, Leonie Moyle, Alessandro Flammini, Esfandiar Haghverdi, and Christopher Raphael all taught excellent classes, and demand my highest respect for their teaching abilities and the high quality research they conduct.

I would like to acknowledge my family. Without both of my parents I would have never made it as far as I have. They have both been equally instrumental in my achievements. My mother, Linda Anderson, who pulled herself up from her bootstraps while raising two kids as a single mom, showed me the importance of hard work and determination.  My father, Clifford Clark, who I resemble in so many ways, has always been there rooting me on; from little league to my dissertation defense. He has always placed the well being of my sister and myself as his top priority. 

Finally, and most importantly, I have to thank my wife, Sergelen Ariunbaatar, who has shown me infinite patience and love. She has always encouraged me to continue when I've become frustrated, and has been extremely understanding of the long hours demanded of a graduate student.   


\end{acknowledgements}
